\section{Introduction}
\label{sec:introduction}

The ionized gas in HII regions is dynamically structured by stellar feedback,
pre-existing molecular cloud turbulence, photoevaporative flows, expanding
shells, and the geometry of the surrounding density field. These processes
inject and redistribute kinetic energy over a range of spatial scales, so the
observed line-of-sight velocity field is not expected to be described by a
single characteristic velocity dispersion. A scale-dependent statistic is
therefore required to separate small-scale velocity fluctuations from
large-scale coherent motions.

In this work we use velocity structure functions (VSFs) measured from LVM
integral-field observations to quantify the spatial scaling of ionized-gas
motions in a sample of Galactic HII regions. The second-order VSF,
\begin{equation}
    S_2(\ell)
    =
    \left\langle
    \left| v(\mathbf{r}+\boldsymbol{\ell}) - v(\mathbf{r}) \right|^2
    \right\rangle ,
\end{equation}
measures the mean squared velocity difference between pairs of spaxels
separated by projected physical distance $\ell$. A power-law dependence
$S_2(\ell)\propto \ell^m$ provides a compact description of the velocity
scaling, while curvature or breaks in $S_2$ indicate that different physical
processes dominate at different scales.

The paper is organized as a progression from morphology to kinematics to
statistical interpretation. We first define the regions and the spatial masks
used for the analysis. We then compute second-order VSFs for all regions and
emission-line tracers, fit those VSFs with continuous segmented power laws, and
identify the spatial scales over which the fitted slope is closest to the
reference value $m=1.66$. Finally, we extend the analysis beyond $S_2$ by
examining the full velocity-difference distribution and comparing higher-order
VSFs against $S_3$.

\section{Data, Tracers, and Spatial Selection}
\label{sec:data-selection}

For each HII region, we use the LVM emission-line measurements of the ionized
gas velocity field. The main tracers considered in the VSF analysis are [NII],
[SII], and [SIII]. These lines sample related but not identical ionization
conditions: [SII] and [NII] preferentially trace lower-ionization gas, while
[SIII] traces more highly ionized material. Comparing their VSFs therefore
allows us to test whether the measured velocity scaling is common to the
ionized region as a whole or varies with ionization structure.

Spaxels are required to satisfy the adopted signal-to-noise and finite-velocity
criteria before entering the VSF calculation. In addition, the main analysis
uses selected subregions defined from the [NII] flux morphology. This selection
is designed to isolate the bright ionized structures associated with each HII
region and to reduce contamination from faint diffuse emission or unrelated
field components. The same spatial selection is then applied consistently to
the different velocity tracers, while retaining the line-specific quality cuts.

\subsection{Morphological Summary of the Selected Regions}
\label{sec:summary-maps}

The first figure provides the morphological reference frame for the rest of the
paper. It shows representative [NII] flux maps for three HII regions, together
with the selected ionized subregions and bright source locations. This figure
establishes where the velocity pairs used in the VSFs are drawn from and shows
how the selected subregions relate to the full LVM field of view.

This step is important because the VSF is sensitive to the spatial domain over
which pairs are drawn. Including unrelated diffuse emission can artificially
mix different kinematic components, while an overly restrictive mask can remove
the large separations needed to identify turnover scales. The selected-subregion
approach is therefore a compromise: it preserves the main bright structures
while keeping enough spatial dynamic range for a scale-dependent analysis.

The corresponding pipeline products are generated in
\texttt{manuscript\_pipeline/01\_summary\_regions/}. The summary figure is
\texttt{Figures/manuscript/01\_summary\_regions/nii\_flux\_maps\_all\_subregions\_with\_sources\_and\_subregions.png}.

\section{Second-Order Velocity Structure Functions}
\label{sec:vsf2}

After defining the spatial masks, we compute $S_2(\ell)$ for each HII region
and each available emission-line tracer. Pair separations are converted to
projected physical distances using the adopted distance to each region and are
binned logarithmically. In each radial bin, the VSF is the mean of
$|\Delta v|^2$ over all valid pairs in that bin. The result is a set of
directly comparable $S_2(\ell)$ curves for [NII], [SII], and [SIII] across the
sample.

The uncertainty in each VSF bin is estimated by bootstrap resampling the input
spaxels rather than by resampling velocity pairs. This distinction matters
because the pair values are highly non-independent: a single spaxel contributes
to many pair separations. For each bootstrap realization, we resample the
spaxels with replacement and recompute the binned VSF. The uncertainty assigned
to a bin is the standard deviation of the bootstrap distribution,
\begin{equation}
    \sigma_{\mathrm{boot}}(\ell_j)
    =
    \mathrm{std}
    \left[
    S_2^{(b)}(\ell_j)
    \right]_{b=1}^{N_{\mathrm{boot}}}.
\end{equation}
These bootstrap errors are used both in the plotted VSF uncertainties and in
the weighted segmented fits described below.

The main $S_2$ figure shows the measured VSFs for all selected subregions.
This figure is the observational anchor for the paper. It shows that the VSF
amplitude, radial curvature, and tracer-to-tracer ordering vary across the
sample. Some regions show relatively smooth power-law-like behavior over much
of the measured range, while others show flattening, steepening, or turnover
features. These differences motivate replacing a single-slope description with
a segmented model.

The selected-subregion VSF figure is
\texttt{Figures/manuscript/02\_vsf2\_multifit/multi\_fits\_vsf\_absolute\_selected\_subregions\_mixed\_thresholds.png}.

\section{Segmented Power-Law Fits}
\label{sec:segmented-fits}

The VSF curves show that a single power law is not always an adequate summary
of the velocity scaling. To quantify this behavior, we fit each $S_2(\ell)$
with continuous segmented power laws in log-log space. For a model with $m$
segments and $m-1$ breakpoints, the fitted relation can be written as
\begin{equation}
    y(x)
    =
    a + b_1 x
    +
    \sum_{k=1}^{m-1} c_k \max(0,x-\tau_k),
\end{equation}
where $x=\log_{10}\ell$, $y=\log_{10}S_2$, and $\tau_k$ are the breakpoint
locations. This parameterization enforces continuity at the breakpoints while
allowing the local slope to change from one segment to the next.

For each region and tracer, we fit candidate models with one, two, and three
segments. At fixed breakpoint locations, the model is linear in the
coefficients and is solved with weighted least squares using the bootstrap VSF
errors. The preferred number of segments is selected with the Bayesian
Information Criterion (BIC). A more complex model is accepted only when the BIC
improvement over the best simpler model exceeds the adopted threshold, so that
breakpoints are retained only when they provide a substantial improvement to
the description of the VSF.

The segmented-fit diagnostic figure shows the measured VSFs and the selected
best-fitting models. Its role is twofold. First, it verifies that the BIC
selection is capturing real curvature or changes in slope rather than fitting
noise. Second, it translates the visually complex VSF curves into a compact set
of quantities: segment slopes, breakpoint scales, and fit uncertainties. These
quantities become the basis for the cross-region comparison in the next
section.

The selected-subregion fit table and diagnostic figure are
\texttt{Figures/manuscript/03\_slopes\_turnover/segmented\_vsf\_breakpoint\_fits\_selected\_subregions\_logflux2.8.csv}
and
\texttt{Figures/manuscript/03\_slopes\_turnover/segmented\_vsf\_bestfits\_selected\_subregions\_logflux2.8.png}.

\section{Breakpoint Scales Associated with the Reference Slope}
\label{sec:slope-turnover}

The segmented fits provide multiple slopes and, in some cases, multiple
breakpoints for a single VSF. To compare all regions in a uniform way, we
select the fitted segment whose slope is closest to $m=1.66$. This value is
used as a reference scaling rather than as a forced model constraint. For the
selected segment, we record the breakpoint that defines the right-hand boundary
of that segment. This gives the largest spatial scale over which the
near-reference-slope behavior persists within the fitted model.

The slope--breakpoint figure places every accepted segmented fit into this
common parameter space. The vertical coordinate is the selected segment slope,
with uncertainty taken from the corresponding fitted slope error. The
horizontal coordinate is the associated breakpoint scale. Points are colored by
emission-line tracer. This representation asks whether the segment closest to
the reference scaling occurs over similar physical scales in different HII
regions, and whether [NII], [SII], and [SIII] occupy systematically different
parts of the diagram.

This figure follows naturally from the segmented-fit diagnostics: instead of
showing each VSF individually, it compresses the fitted models into a
population-level comparison. If the points cluster around a common breakpoint
scale, that would suggest a characteristic transition scale in the ionized gas.
If the scales differ strongly by tracer or by region, the velocity scaling is
more likely controlled by local feedback geometry, ionization stratification,
or region-specific structure.

The current slope--breakpoint figure is
\texttt{Figures/manuscript/03\_slopes\_turnover/slope\_closest\_1p66\_vs\_breakpoint\_scale\_selected\_subregions\_logflux2.8.png}.

\section{Velocity-Difference Distributions in the Flame Nebula}
\label{sec:flame-3d}

The previous sections reduce each velocity field to $S_2(\ell)$ and then to a
small number of fitted slopes and breakpoints. Before moving to higher-order
VSFs for the full sample, we examine one region in more detail to show what
information is compressed into the VSF. For the Flame Nebula [NII] velocity
field, we visualize the velocity-difference distribution as a function of
separation, together with statistics derived from that distribution.

This figure emphasizes that changes in $S_2(\ell)$ are not only changes in a
mean squared velocity difference. They can also reflect changes in the shape,
width, and tails of the pairwise velocity-difference distribution. In
particular, broad or non-Gaussian wings at specific separations can influence
higher-order moments more strongly than $S_2$. The Flame Nebula visualization
therefore provides a bridge between the second-order segmented-fit analysis and
the higher-order VSF analysis that follows.

The corresponding figure is
\texttt{Figures/manuscript/04\_flame\_kurtosis\_3d/Flame\_NII\_refactored\_3D.png}.

\section{Higher-Order VSFs Relative to $S_3$}
\label{sec:higher-order-vsfs}

The second-order VSF measures variance-like velocity structure, but it does not
fully characterize the velocity-difference distribution. To test whether the
ionized-gas motions are self-similar across moments, we compute VSFs of
multiple orders and compare them to the third-order VSF. For order $p$,
\begin{equation}
    S_p(\ell)
    =
    \left\langle
    \left| \Delta v(\ell) \right|^p
    \right\rangle .
\end{equation}
The relation between $S_p$ and $S_3$ provides a relative scaling measurement
that is less dependent on the absolute separation axis and is commonly used to
diagnose departures from simple self-similar behavior.

The VSF-order figure shows the relation between $S_p$ and $S_3$ for
$p=1$--$5$. If the velocity field followed a single self-similar hierarchy, the
slopes of these relations would follow a predictable monotonic trend with
order. Deviations from that trend indicate intermittency or changes in the
relative importance of rare, large velocity differences. This section therefore
extends the analysis from the scale dependence of $S_2$ to the moment
dependence of the full pairwise velocity field.

The main higher-order VSF figure is
\texttt{Figures/manuscript/06\_vsf\_orders/VSF\_orders1to5\_vs\_order3\_allregions.png}.

\section{Final Slope--Order Comparison by Region}
\label{sec:slope-order-final}

The final figure condenses the higher-order VSF measurements into a
region-by-region comparison of slope versus order. For each HII region, the
fitted slope of $S_p$ relative to $S_3$ is plotted as a function of $p$. This
representation makes it possible to compare the intermittency behavior of
different regions directly. Regions with similar curves have similar relative
moment scaling, whereas regions with steeper or flatter slope--order relations
show differences in the tails of the velocity-difference distribution.

This final step connects back to the segmented $S_2$ analysis. The
slope--breakpoint figure identifies the scale over which a reference-like
second-order slope persists, while the slope--order figure asks whether the
same regions also share a common hierarchy of higher-order moments. Agreement
between these diagnostics would suggest a coherent turbulent scaling regime;
disagreement would indicate that similar $S_2$ slopes can arise from different
underlying velocity-difference distributions.

The final figure is
\texttt{Figures/manuscript/07\_final\_slope\_vs\_order/vsf\_vs3\_slope\_vs\_order\_by\_region.png}.

\section{Discussion}
\label{sec:discussion}

The analysis shows that the ionized-gas velocity fields in the HII-region
sample are not generally described by a single scale-free power law. The
measured $S_2(\ell)$ curves often contain curvature or slope changes that are
captured by the segmented fits. In physical terms, this means that the
line-of-sight velocity differences do not grow with separation at a single rate
over the full observed range. Instead, the data favor regimes in which the
dominant kinematic structure changes with scale.

The breakpoints should not be interpreted automatically as sharp physical
boundaries. They are empirical transition scales in a projected statistic. A
breakpoint may correspond to the characteristic size of a coherent shell, the
width of an ionized ridge, the scale over which photoevaporative flows remain
correlated, or the transition from local feedback-driven motions to larger
cloud-scale velocity gradients. Projection effects can also shift or broaden
these transitions. The strength of the segmented-fit approach is therefore not
that it assigns a unique physical mechanism to each break, but that it provides
a reproducible way to compare the scale dependence of the velocity field across
regions and tracers.

The tracer comparison is central to this interpretation. [NII], [SII], and
[SIII] do not necessarily sample the same gas parcels or ionization zones.
If their VSFs have similar slopes and breakpoints within a region, the dominant
kinematic structure is likely shared across the ionized gas. Conversely,
systematic offsets between tracers can indicate ionization stratification,
differential coupling to expanding structures, or different sensitivity to
dense interfaces and diffuse interiors. The slope--breakpoint diagram is useful
because it separates these effects into two questions: whether a tracer reaches
a reference-like slope, and over what spatial scale that behavior persists.

The reference slope of $m=1.66$ is used here as a comparative anchor rather
than as a claim that all regions obey a universal turbulent law. Selecting the
segment closest to this value allows us to ask where each VSF most closely
resembles that scaling and where it departs from it. The associated right-hand
breakpoint marks the largest fitted scale for that segment. Regions whose
near-reference segments end at similar scales may share a common driving or
decorrelation scale; regions with very different breakpoints likely reflect
differences in size, feedback geometry, or environmental structure.

The Flame Nebula velocity-difference visualization and the higher-order VSFs
show why the second-order analysis is only part of the story. Two regions can
have similar $S_2$ slopes while having different velocity-difference tails.
Higher-order moments are more sensitive to rare, large velocity differences,
and therefore probe intermittency that is mostly hidden in $S_2$. If the
slope--order curves differ strongly among regions, then the ionized-gas
turbulence is not fully described by the second-order turnover scale alone.
Instead, the results point to a combination of scale-dependent energy injection
and region-dependent intermittency.

Several limitations should be kept in mind. The analysis uses projected
separations and line-of-sight velocities, so the measured VSFs are projected
statistics of a three-dimensional flow. The selected-subregion masks reduce
contamination from unrelated emission but also define the maximum available
separation and can influence the large-scale VSF behavior. The bootstrap errors
capture sampling variance and pair non-independence, but they do not remove
possible systematics from velocity fitting, distance uncertainties, or
line-of-sight blending. These limitations affect the absolute interpretation of
individual breakpoints, but they do not erase the comparative value of applying
the same measurement and fitting procedure to every region.

\section{Conclusions}
\label{sec:conclusions}

We have developed a VSF-based framework for comparing the scale-dependent
kinematics of ionized gas in Galactic HII regions. The main conclusions are:
\begin{enumerate}
    \item We compute second-order VSFs for [NII], [SII], and [SIII] in selected
    ionized subregions, using bootstrap resampling of spaxels to estimate the
    uncertainties in each radial bin.
    \item The measured $S_2(\ell)$ curves frequently show curvature or changes
    in slope, motivating continuous segmented power-law fits rather than a
    single-slope description.
    \item BIC-selected segmented models provide empirical breakpoint scales and
    segment slopes for each region and tracer. These quantities summarize where
    the velocity scaling changes across the observed spatial range.
    \item By selecting the segment whose slope is closest to $m=1.66$ and
    recording the breakpoint at its right-hand boundary, we compare the spatial
    extent of near-reference scaling across regions and emission-line tracers.
    \item The Flame Nebula velocity-difference distribution and the higher-order
    VSFs show that the second-order VSF is a compressed statistic of a richer,
    sometimes intermittent velocity field.
    \item The final slope--order comparison demonstrates that region-to-region
    differences persist beyond $S_2$, indicating that both the scale dependence
    and the moment dependence of the velocity field are needed to describe the
    ionized-gas kinematics.
\end{enumerate}

Together, these results show that LVM velocity maps can be used to measure not
only the amplitude of ionized-gas motions, but also the characteristic scales
and higher-order structure of those motions across a diverse sample of HII
regions.
